\section{Prompts zur Erstellung der Abbildungen}
\label{anhang:prompst_abbildungen}

\textbf{Abbildung 1: Systemablauf innerhalb der Webanwendung}
\newline
Prompt: 
\enquote{Erstelle ein kompaktes, Prozessdiagramm im Querformat, das sich gut zwischen Text auf einer DINA4 Seite einfügen lässt. Verwende ausschließlich Schwarz, Weiß und helle Grautöne, einen weißen Hintergrund.

Die Prozessschritte lauten:\newline
1 BROWSER: Aufgabe eingeben\newline
2 API: Eingabe prüfen\newline
3 BACKEND: Prompt bilden\newline
4 OPENAI: Testfälle erzeugen\newline
5 API: Antwort zurückgeben\newline
6 BROWSER: JSON darstellen

Verbinde die Boxen mit n Pfeilen. Beschrifte die Verbindungen mit:\newline
1 -> 2: POST /api\newline
2 -> 3: validiert\newline
3 -> 4: HTTPS\newline
4 -> 5: output\_text\newline
5 -> 6: JSON

Platziere sämtliche Pfeilbeschriftungen mit ausreichend Abstand zu den Boxen und Pfeilen. Die Beschriftungen sollen mit keinem anderen Elementen überlappen.}

\textbf{Abbildung 2: Architektur der Webanwendung}
\newline
Prompt:
\enquote{Erstelle ein Architekturdiagramm im gleichen Format und Style wie das Prozessdiagramm.

Stelle drei Komponenten nebeneinander dar:

1. Webbrowser als Client: \newline
- Zeigt das über die GET / Route das UI mit den Dateien index.html, styles.css und app.js vom Server bereitgestellt wird.\newline
- Zeigt außerdem das über die POST /api Route die Anfrage an das Backend geschickt wird.

2. FastAPI Backend als Server:\newline
-  GET / worüber die statischen Dateien an den Client geschickt werden\newline
- POST /api wo die Anfragen ankommen, geprüft werden, der System-Promt eingebaut wird und die Anfrage an die OpenAI API geschickt wird

3. OpenAI API:\newline
Verarbeitet die Anfrage des Backends, erstellt die Testfälle und gibt sie zurück}